
\chapter{Examples and Applications}

\section{Examples of Small Groups}
\begin{exmp}[The cyclic group \(C_n\)]
Let \(G = \langle g \mid g^n = e \rangle\). Over \(\mathbb{C}\), the irreducible representations are 1-dimensional, given by
\[
\rho_k(g) = \zeta^k, \qquad k = 0,1,\dots,n-1,
\]
where \(\zeta\) is a primitive \(n\)-th root of unity. The character table is simply \(\chi_k(g^m) = \zeta^{km}\).
\end{exmp}

\begin{exmp}[The symmetric group \(S_3\)]
\(S_3\) has three conjugacy classes:
\[
C_1 = \{1\}, \quad C_2 = \{(12), (13), (23)\}, \quad C_3 = \{(123), (132)\}.
\]
Their sizes are \(1, 3, 2\). The character table is:

\[
\begin{array}{c|ccc}
 & 1 & (12) & (123) \\
\hline
\chi_1 & 1 & 1 & 1 \\
\chi_2 & 1 & -1 & 1 \\
\chi_3 & 2 & 0 & -1
\end{array}
\]

Check of orthogonality:
\[
\langle \chi_3, \chi_3 \rangle = \frac{1}{6}(2^2 + 3 \cdot 0 + 2 \cdot (-1)^2) = \frac{4 + 2}{6} = 1.
\]
Check of dimensions:
\[
1^2 + 1^2 + 2^2 = 6 = |S_3|.
\]
\end{exmp}

\begin{exmp}[The alternating group \(A_4\)]

\(A_4\) has four conjugacy classes:
\[
C_1 = \{1\}, \quad C_2 = \{(12)(34), (13)(24), (14)(23)\}, \quad C_3 = \{(123), (142), (134), (243)\}, \quad C_4 = \{(132), (124), (143), (234)\}.
\]
Their sizes are \(1, 3, 4, 4\). Let \(\omega = e^{2\pi i / 3}\). The character table is:

\[
\begin{array}{c|cccc}
 & 1 & (12)(34) & (123) & (132) \\
\hline
\chi_1 & 1 & 1 & 1 & 1 \\
\chi_2 & 1 & 1 & \omega & \omega^2 \\
\chi_3 & 1 & 1 & \omega^2 & \omega \\
\chi_4 & 3 & -1 & 0 & 0
\end{array}
\]

How to obtain \(\chi_4\): It is the standard 3-dimensional representation of \(S_4\) restricted to \(A_4\), or equivalently \(\operatorname{Ind}_{C_3}^{A_4} \chi\) where \(\chi\) is a nontrivial character of \(C_3\). Check:
\[
\langle \chi_4, \chi_4 \rangle = \frac{1}{12}(3^2 + 3 \cdot (-1)^2 + 4 \cdot 0 + 4 \cdot 0) = \frac{9 + 3}{12} = 1.
\]
Check of dimensions:
\[
1^2 + 1^2 + 1^2 + 3^2 = 12 = |A_4|.
\]
\end{exmp}
\begin{remark}
The character table of \(A_4\) demonstrates a key principle: irreducible representations of a quotient group (here \(A_4 / V_4 \cong C_3\)) lift to representations of the whole group. The three 1-dimensional characters come from this lifting.
\end{remark}
\section{Application to the Classification of Finite Groups}

One of the crowning achievements of twentieth-century mathematics is the 
\textbf{classification of finite simple groups} (CFSG). The theorem states 
that every finite simple group belongs to one of the following families:

\begin{enumerate}
    \item cyclic groups of prime order $C_p$;
    \item alternating groups $A_n$ for $n \ge 5$;
    \item groups of Lie type (e.g., $\mathrm{PSL}_n(q)$, unitary groups, 
        symplectic groups, and their twisted variants);
    \item the $26$ sporadic simple groups, including the Monster, the 
        Baby Monster, and Mathieu groups.
\end{enumerate}

The proof of this monumental theorem spans thousands of pages and relies on 
deep interactions between group theory, geometry, and number theory. 
\textbf{Representation theory} — especially character theory and modular 
representation theory — plays an indispensable role throughout the 
classification programme. The following points 
illustrate the most direct connections.

\medskip

\noindent \textbf{Detecting normal subgroups via characters.} 
A character $\chi$ of a finite group $G$ is called \emph{faithful} if its 
kernel is trivial. More generally, the intersection of the kernels of all 
irreducible characters is the trivial subgroup. Hence, if we can compute 
the character table of $G$, we can read off all normal subgroups: they are 
exactly intersections of kernels of irreducible characters. Thus, to show 
that a group is simple, it suffices to verify that no non‑trivial character 
has a non‑trivial kernel. In practice, this criterion is used to rule out 
possible normal subgroups in candidate groups.

\medskip

\noindent \textbf{Modular representations and local analysis.} 
The classification proof heavily uses \emph{local subgroups} — normalisers 
of non‑trivial $p$-subgroups. To understand their structure, one needs 
representation theory over fields of characteristic $p$. Brauer's theory of 
modular characters (or $p$-modular characters) provides a powerful tool to 
study the $p$-local structure of a group. For instance, Brauer's 
\emph{first main theorem} establishes a bijection between blocks of a group 
and blocks of its $p$-local subgroups, which is crucial for inductive 
arguments. Without modular representation theory, the analysis of 
$p$-subgroups and their normalisers would be intractable.

\medskip

\noindent \textbf{Induced representations and the Clifford–Mackey machinery.} 
When constructing new simple groups or proving that a given group is 
simple, one often needs to understand how representations behave under 
induction from subgroups. Mackey's irreducibility criterion and Clifford 
theory allow one to determine when an induced representation remains 
irreducible. This is used, for example, to build the so‐called 
\emph{complex reflection groups} and to distinguish different Lie‑type 
groups from each other.

\medskip

\noindent \textbf{Galois conjugacy and Schur indices.} 
The character table of a group has entries in cyclotomic fields. The action 
of the Galois group on these entries imposes strong constraints on the 
possible group structure. Moreover, the Schur index — the smallest field 
over which a representation can be realised — can obstruct the existence of 
certain subgroups or extensions. These arithmetic considerations were 
essential in eliminating many potential simple groups during the 
classification.

\medskip

\noindent \textbf{The role of the Monster and vertex operator algebras.} 
The largest sporadic simple group, the Monster, has a 196‑883‑dimensional 
irreducible representation over $\mathbb{C}$. Its character table, computed 
using representation‑theoretic techniques, revealed surprising connections 
with modular functions (the $j$-invariant) and led to monstrous moonshine. 
This interplay between representation theory and number theory not only 
confirmed the existence of the Monster but also opened new avenues in 
mathematical physics.

\medskip

For a beginning student, the key takeaway is this: \emph{representation 
theory provides a linear lens through which the hidden structure of a 
finite group becomes visible.} It transforms questions about subgroups and 
normality into questions about matrices, characters, and algebraic 
equations. While the full classification theorem is far beyond the scope of 
this course, understanding the role of characters and induced 
representations gives you a glimpse into the power of the subject — a power 
that ultimately helped mathematicians map out all finite building blocks of 
symmetry.

\bigskip

\noindent \textbf{Further reading.} For an accessible introduction, see \cite{alperin1995groups} and \cite{aschbacher2004status}.
\section{Applications to Number Theory:Artin L-functions}
In number theory, Dirichlet characters are precisely the irreducible characters of the finite abelian group \((\mathbb{Z}/m\mathbb{Z})^\times\). More generally, Artin L-functions are attached to representations of Galois groups; their study relies heavily on character theory.


In this subsection, let \(L/K\) be a finite Galois extension of number fields, with Galois group \(G := \mathrm{Gal}(L/K)\). Let \(\mathcal{O}_K\) and \(\mathcal{O}_L\) be their rings of integers (see Definition \ref{def:ring of integers}). 

\subsection{Definition and Basic Properties of Artin L-functions}


\begin{defn}
The \textbf{absolute norm} of  a nonzero prime ideal $\mathfrak p\subset\mathcal O_K$ is 
\[
{N(\mathfrak p) := |\mathcal O_K/\mathfrak p|.}
\]
\end{defn}

We first recall the decomposition of \(\mathfrak{p}\) in \(\mathcal{O}_L\).
\begin{defn}
By Theorem \ref{thm:to prime ideals}(replace $K$ with $L$), there exist prime ideals \(\mathfrak{P}_1, \dots, \mathfrak{P}_g \subset \mathcal{O}_L\) and positive integers \(e_1, \dots, e_g\) such that
\[
\mathfrak{p}\mathcal{O}_L = \mathfrak{P}_1^{e_1} \cdots \mathfrak{P}_g^{e_g}.
\]
The integer \(e_i\) is called the \textbf{ramification index} of \(\mathfrak{P}_i\) over \(\mathfrak{p}\).

We say that \(\mathfrak{p}\) is \textbf{unramified} in \(L/K\) if
\[
e_i = 1 ,\quad \forall i \in\{1, \dots, g\}.
\]
\end{defn}


Now assume the nonzero prime ideal $\mathfrak p\subset\mathcal O_K$ is unramified in \(L/K\). Let \(\mathfrak{P} \subset \mathcal{O}_L\) be a prime ideal such that \(\mathfrak{P} \cap \mathcal{O}_K = \mathfrak{p}\) (i.e.,\(\mathfrak{P}\) is a prime ideal lying over \(\mathfrak{p}\)). 



\begin{lemma}
There exists a unique element $\sigma\in G$ such that
\[
\sigma(\mathfrak{P}) = \mathfrak{P}
\]
and
\[
\sigma(\alpha) \equiv \alpha^{N(\mathfrak{p})} \pmod{\mathfrak{P}},
\quad \forall \alpha \in \mathcal{O}_L.
\]

The \textbf{Frobenius element} at $\mathfrak{P}$ is defined by
\[
\mathrm{Frob}_{\mathfrak{P}} := \sigma.
\]

Moreover, if \(\mathfrak{P}'\) is another prime ideal lying over \(\mathfrak{p}\), then \(\mathfrak{P}' = g(\mathfrak{P})\) for some \(g \in G\), and one has
\[
\mathrm{Frob}_{\mathfrak{P}'} = g \,\mathrm{Frob}_{\mathfrak{P}}\, g^{-1}.
\]
Thus the conjugacy class of \(\mathrm{Frob}_{\mathfrak{P}}\) is independent of the choice of \(\mathfrak{P}\). This conjugacy class is called \textbf{Frobenius conjugacy class} at $\mathfrak{p}$, denoted by $\mathrm{Frob}_{\mathfrak{p}}$.
\end{lemma}

\begin{proof}
Since \(L/K\) be a finite Galois extension, we can choose \(\theta \in \mathcal{O}_L\) such that \(L = K(\theta)\). Let \(f(x) \in \mathcal{O}_K[x]\) be minimal polynomial of $\theta$ over \(K\). It splits completely in \(L[x]\),
\[
f(x) = (x - \theta_1)(x - \theta_2)\cdots(x - \theta_n),
\]
with \(\theta_1 = \theta\) and all \(\theta_i \in \mathcal{O}_L\).

Reducing modulo \(\mathfrak{P}\), we get
\[
\bar{f}(x) = (x - \bar{\theta}_1)(x - \bar{\theta}_2)\cdots(x - \bar{\theta}_n),
\]
where \(\bar{\theta}_i = \theta_i + \mathfrak{P}\). Since \(\mathfrak{p}\) is unramified, \(\bar{f}\) has no repeated roots; hence the \(\bar{\theta}_i\)'s are all distinct.

In the residue field $\mathcal O_L/\mathfrak{P}$, the Frobenius automorphism sends \(\bar{\theta}_1\) to \(\bar{\theta}_1^{N(\mathfrak{p})}\), which is again a root of \(\bar{f}\). Therefore there exists a unique \(j\) such that
\[
\bar{\theta}_1^{N(\mathfrak{p})} = \bar{\theta}_j,
\]
or equivalently,
\[
\theta_j \equiv \theta^{N(\mathfrak{p})} \pmod{\mathfrak{P}}.
\]

Define
\[
\sigma:K\to K,\quad \theta\mapsto\theta_j,
\]

Since \(L = K(\theta)\), every element \(\alpha \in L\) can be written as
\[
\alpha = \frac{f(\theta)}{h(\theta)}, \qquad (f(x), h(x)) \in K[x]^2, \quad h(\theta) \neq 0.
\]
Define \(\sigma: L \to L\) by
\[
\sigma\!\left(\frac{f(\theta)}{h(\theta)}\right) := \frac{f(\theta_j)}{h(\theta_j)}.
\]
This is well-defined: if \(\frac{f_1(\theta)}{h_1(\theta)} = \frac{f_2(\theta)}{h_2(\theta)}\), then
\[
f_1(x)h_2(x) - f_2(x)h_1(x)
\]
is divisible by the minimal polynomial of \(\theta\), hence vanishes at \(x = \theta_j\), so the two rational expressions in \(\theta_j\) are equal. It is straightforward to check that \(\sigma\) preserves addition and multiplication, and fixes \(K\). Hence \(\sigma \in \mathrm{Gal}(L/K)\). It is easy to check $\sigma(\mathcal{O}_L)\subseteq \mathcal{O}_L$.

For any \(\alpha = h(\theta) \in \mathcal{O}_L\), we have
\[
\sigma(\alpha) = h(\theta_j) \equiv h(\theta^{N(\mathfrak{p})}) \equiv h(\theta)^{N(\mathfrak{p})} = \alpha^{N(\mathfrak{p})} \pmod{\mathfrak{P}},
\]

We now prove that \(\sigma(\mathfrak{P}) = \mathfrak{P}\).  
Take any \(\beta \in \mathfrak{P}\). From the congruence
\[
\sigma(\beta) \equiv \beta^{N(\mathfrak{p})} \equiv0\pmod{\mathfrak{P}},
\]
so \(\sigma(\mathfrak{P}) \subseteq \mathfrak{P}\).

Since \(\sigma\) is a bijection on \(\mathcal{O}_L\), $|\sigma(\mathfrak{P})| = |\mathfrak{P}|$, then
\[
\sigma(\mathfrak{P}) = \mathfrak{P}.
\]


Uniqueness follows from the fact that an element of \(G\) is uniquely determined by its action on \(\theta\), and the condition \(\sigma(\theta) \equiv \theta^{N(\mathfrak{p})} \pmod{\mathfrak{P}}\) singles out a unique conjugate root \(\theta_j\) among the distinct roots.


Let \(\mathfrak{P}'\) be another prime ideal above \(\mathfrak{p}\). Since \(G\) acts transitively on the set of prime ideals above \(\mathfrak{p}\), there exists \(g \in G\) such that
\[
\mathfrak{P}' = g(\mathfrak{P}).
\]

Let \(\sigma = \mathrm{Frob}_{\mathfrak{P}}\). First,
\[
(g\sigma g^{-1})(\mathfrak{P}') = g\sigma(g^{-1}(g(\mathfrak{P}))) = g\sigma(\mathfrak{P}) = g(\mathfrak{P}) = \mathfrak{P}'.
\]
Second, for any \(\alpha \in \mathcal{O}_L\), since \(g^{-1}(\alpha) \in \mathcal{O}_L\), we have
\[
\sigma(g^{-1}(\alpha)) \equiv g^{-1}(\alpha)^{N(\mathfrak{p})} \pmod{\mathfrak{P}}.
\]
Applying \(g\) to both sides gives
\[
g\sigma(g^{-1}(\alpha)) \equiv \alpha^{N(\mathfrak{p})} \pmod{g(\mathfrak{P})} = \pmod{\mathfrak{P}'}.
\]
Thus \(g\sigma g^{-1}\) satisfies the defining properties of \(\mathrm{Frob}_{\mathfrak{P}'}\). By uniqueness, we obtain the desired equality\[
\mathrm{Frob}_{\mathfrak{P}'} = g \,\mathrm{Frob}_{\mathfrak{P}}\, g^{-1}.
\]
\end{proof}

For convenience, we first define\textbf{ the  unramified case}.

\begin{defn}[Artin L-function, the  unramified case]
Let \((\rho,V)\) be a representation of \(G\) with character \(\chi_\rho\). The \textbf{Artin L-function} \(L(s, \chi_\rho, L/K)\) is defined as an Euler product
\[
L(s, \chi_\rho, L/K) = \prod_{\mathfrak{p}} \det \left( I - \rho(\mathrm{Frob}_{\mathfrak{p}}) N(\mathfrak{p})^{-s} \right)^{-1},
\]
where the product is over each prime ideal \(\mathfrak{p}\) of \(\mathcal{O}_K\) that is unramified in \(L/K\), and $\rho(\mathrm{Frob}_{\mathfrak{p}})$ is the conjugacy class of  \(\rho(\mathrm{Frob}_{\mathfrak{P}})\). The product is well-defined since the value of $\det \left( I - \rho(\mathrm{Frob}_{\mathfrak{P}}) N(\mathfrak{p})^{-s} \right)^{-1}$does not depend on the choice of $\mathfrak{P}$.

$L_{\mathfrak{p}}(s, \chi_\rho, L/K):= \det \left( I - \rho(\mathrm{Frob}_{\mathfrak{p}}) N(\mathfrak{p})^{-s} \right)^{-1}$ is called the local factor of \(L(s, \chi_\rho, L/K)\) at \(\mathfrak{p}\).

Sometimes we write $L(s, \chi_\rho, L/K)$ simply as $L(s, \rho, L/K)$
\end{defn}

Here are some basic properties of the Artin L-function.
\begin{prop}
Let \(\rho_1, \rho_2\) be two complex representations of \(G\). Then
\[
L(s, \chi_{\rho_1 \oplus \rho_2}, L/K) = L(s,\chi_{ \rho_1}, L/K) \cdot L(s,\chi_{ \rho_2}, L/K).
\]
\end{prop}

\begin{proof}
By definition,
\[
L(s,\chi_ \rho, L/K) = \prod_{\mathfrak{p}} \det\!\left( I - \rho(\mathrm{Frob}_{\mathfrak{p}}) N(\mathfrak{p})^{-s} \right)^{-1},
\]
where the product is over the unramified prime ideals \(\mathfrak{p}\) of \(\mathcal{O}_K\).

For each such \(\mathfrak{p}\), since \(\rho = \rho_1 \oplus \rho_2\), the matrix \(\rho(\mathrm{Frob}_{\mathfrak{p}})\) is block-diagonal:
\[
\rho(\mathrm{Frob}_{\mathfrak{p}}) =
\begin{pmatrix}
\rho_1(\mathrm{Frob}_{\mathfrak{p}}) & 0 \\
0 & \rho_2(\mathrm{Frob}_{\mathfrak{p}})
\end{pmatrix}.
\]
Therefore,
\[
I - \rho(\mathrm{Frob}_{\mathfrak{p}}) N(\mathfrak{p})^{-s}
=
\begin{pmatrix}
I - \rho_1(\mathrm{Frob}_{\mathfrak{p}}) N(\mathfrak{p})^{-s} & 0 \\
0 & I - \rho_2(\mathrm{Frob}_{\mathfrak{p}}) N(\mathfrak{p})^{-s}
\end{pmatrix}.
\]
Using the fact that the determinant of a block-diagonal matrix is the product of the determinants of its blocks, we get
\[
\det\!\left( I - \rho(\mathrm{Frob}_{\mathfrak{p}}) N(\mathfrak{p})^{-s} \right)
=
\det\!\left( I - \rho_1(\mathrm{Frob}_{\mathfrak{p}}) N(\mathfrak{p})^{-s} \right)
\cdot
\det\!\left( I - \rho_2(\mathrm{Frob}_{\mathfrak{p}}) N(\mathfrak{p})^{-s} \right).
\]
Taking inverses and multiplying over all \(\mathfrak{p}\), we obtain
\[
L(s, \chi_{\rho_1 \oplus \rho_2}, L/K)
=
\prod_{\mathfrak{p}}
\det\!\left( I - \rho_1(\mathrm{Frob}_{\mathfrak{p}}) N(\mathfrak{p})^{-s} \right)^{-1}
\cdot
\prod_{\mathfrak{p}}
\det\!\left( I - \rho_2(\mathrm{Frob}_{\mathfrak{p}}) N(\mathfrak{p})^{-s} \right)^{-1}
=
 L(s,\chi_{ \rho_1}, L/K) \cdot L(s,\chi_{ \rho_2}, L/K).
\]

\end{proof}


\begin{prop}[Artin Induction Property]
Let \(H \le G\), and let \(\psi\) be a character of \(H\). Let \(F = L^H\) be the fixed field of \(H\) in \(L\). Then
\[
L(s, \operatorname{Ind}_{H}^{G} \psi, L/K) = L(s, \psi, L/F).
\]
\end{prop}

\begin{proof}
It suffices to compare the local factors at each prime \(\mathfrak{p}\) of \(\mathcal{O}_K\).


Let \(\mathfrak{p}\) be a prime ideal of \(\mathcal{O}_K\) that is unramified in \(L/K\). Let \(\mathfrak{P}\) be a prime ideal of \(\mathcal{O}_L\) lying above \(\mathfrak{p}\). Let \(\sigma = \mathrm{Frob}_{\mathfrak{P}} \in G\) be the Frobenius element at \(\mathfrak{P}\).


For the induced representation \(\operatorname{Ind}_{H}^{G} \psi\), the character at \(\sigma\) is given by the formula
\[
\operatorname{Ind}_{H}^{G} \psi(\sigma) = \sum_{\substack{gH \in G/H \\ g^{-1}\sigma g \in H}} \psi(g^{-1}\sigma g).
\]
Equivalently, the matrix of \(\operatorname{Ind}_{H}^{G} \psi(\sigma)\) is block-diagonal, with blocks corresponding to the orbits of \(\langle \sigma \rangle\) on the coset space \(G/H\). The determinant of \(I - \operatorname{Ind}_{H}^{G} \psi(\sigma) N(\mathfrak{p})^{-s}\) factors over these blocks.

The cosets \(gH\) correspond to the prime ideals \(\mathfrak{P}_g := g(\mathfrak{P})\) of \(\mathcal{O}_L\) lying above \(\mathfrak{p}\). The Frobenius element \(\sigma\) acts on the set of these primes. The cycles of \(\sigma\) on this set correspond exactly to the prime ideals of \(\mathcal{O}_F = \mathcal{O}_{L^H}\) lying below \(\mathfrak{p}\).

For each cycle, say of length \(d\), the block in the determinant contributes
\[
\det\!\left( I - \psi(\sigma^d) N(\mathfrak{p})^{-ds} \right),
\]
which is exactly the local factor of \(L(s, \psi, L/F)\) at the prime of \(\mathcal{O}_F\) lying below the corresponding cycle.


Therefore, after taking the product over all \(\mathfrak{p}\), the local factors on both sides match exactly. Hence the global L-functions are equal:
\[
L(s, \operatorname{Ind}_{H}^{G} \psi, L/K) = L(s, \psi, L/F).
\]
\end{proof}



\subsection{Examples of L-functions}
Some reader may ask: what is L-function? Now we introduce it simply.
\begin{defn}
Let \(K\) be a number field. A local L-factor at a prime ideal \(\mathfrak{p} \subset \mathcal{O}_K\) is typically a rational function in \(N(\mathfrak{p})^{-s}\). The global L-function is the Euler product over all primes:
\[
L(s, \text{data}) = \prod_{\mathfrak{p}} L_{\mathfrak{p}}(s, \text{data})^{-1}.
\]
\end{defn}
Here are some examples.

\begin{itemize}
    \item[(a)] \textbf{Riemann zeta function:}
    \[
    \zeta(s) = \prod_{p\text{ is prime}} \left(1 - p^{-s}\right)^{-1}, \qquad \Re(s) > 1.
    \]

    \item[(b)] \textbf{Dirichlet L-function:}
    \[
    L(s, \chi) = \prod_{p\text{ is prime}} \left(1 - \chi(p) p^{-s}\right)^{-1}.
    \]

    \item[(c)] \textbf{Artin L-function :}
    \[
    L(s, \chi_\rho, L/K) = \prod_{\mathfrak{p} \text{ unramified}} 
    \det\left(1 - \rho(\mathrm{Frob}_{\mathfrak{p}}) N(\mathfrak{p})^{-s}\right)^{-1}.
    \]

    \item[(d)]\textbf{The Hasse--Weil L-function:}
    Let \(E\) be an elliptic curve defined over \(\mathbb{Q}\). For each prime \(p\), let \(\overline{E}_p\) be the reduction of \(E\) modulo \(p\). Define
\[
a_p(E) := p + 1 - \#E(\mathbb{F}_p),
\]
where \(\#E(\mathbb{F}_p)\) is the number of points on the reduced curve over the finite field \(\mathbb{F}_p\). The local L-factor at \(p\) is:
\[
L_p(E, s) :=
\begin{cases}
\left(1 - a_p p^{-s} + p^{1-2s}\right)^{-1}, & \text{if } E \text{ has good reduction at } p,\\[4pt]
\left(1 - a_p p^{-s}\right)^{-1}, & \text{if } E \text{ has split multiplicative reduction},\\[4pt]
\left(1 + a_p p^{-s}\right)^{-1}, & \text{if } E \text{ has non-split multiplicative reduction},\\[4pt]
1, & \text{if } E \text{ has additive reduction}.
\end{cases}
\]
The global Hasse--Weil L-function is then defined as the Euler product
\[
L(E, s) := \prod_{p} L_p(E, s),
\]
which converges for \(\Re(s) > \tfrac{3}{2}\).

\end{itemize}

The Hasse--Weil L-function is the fundamental analytic object attached to an algebraic variety over a number field. It generalizes the Artin L-function from representations of Galois groups to the \'etale cohomology of varieties. 

For the interested reader, you can read \cite{marcus1977number} and \cite{milne2020cft} for further study.



\subsection{The Artin Conjecture}


We first prove the following standard lemmas, which related Brauer's Theorem we studied in Section 9.

\begin{thm}[Clifford]\label{Clifford}
Let \(P\) be a finite group and \(N \trianglelefteq P\). Let $(\rho,V)$ be an irreducible complex representation of \(P\). 

Suppose \((\mu, W) \) is an irreducible subrepresentation of the restricted representation $(\rho|_N,V)$. Define the \textbf{inertia subgroup} of $\mu$ in $P$ as
\[
I := \{ g \in P \mid \mu^g \cong \mu \},
\]
where $\mu^g(n)=\mu(g^{-1}ng),\quad\forall n\in N$.

Then there exists an irreducible representation 
\[
\tau: I \to \mathrm{GL}(U)
\]
such that 
\[
\rho \cong \operatorname{Ind}_{I}^{P} \tau.
\]
\end{thm}

\begin{proof}
By Maschke's Theorem, the restriction \(\rho|_N\) is completely reducible.
Therefore, \(\rho|_N\) decomposes into a direct sum of irreducible \(N\)-invariant  representations, one of which is \(\mu\).

Consider the direct sum of \emph{all} irreducible subrepresentations appearing in \(\rho|_N\) that are isomorphic to \(\mu\). 
Let \(Z \subseteq V\) be the subspace consisting of all vectors that belong to irreducible \(N\)-invariant  subrepresentations isomorphic to \(\mu\).
Since $I = \{ g \in P \mid \mu^g \cong \mu \}$, the subspace \(Z\) is invariant under the action of every element of \(I\). 

Hence, restricting the original matrices of \(\rho\) to the subspace \(Z\) and to the subgroup \(I\) gives a well-defined subrepresentation of \((\rho|_I,V)\) :
\[
\tau: I \to \mathrm{GL}(Z).
\]
Restrict \(\tau\) to \(N\). \(\mu\) appears as one of  direct summands of $\tau$, i.e.,  \(\mu\) is a subrepresentation of \(\tau|_N\).




Since \(\tau\) is a subrepresentation of \(\rho|_I\), the inclusion 
\[
\tau \hookrightarrow \rho|_I
\]
is a non-zero element of \(\operatorname{Hom}_{I}(\tau, \rho|_I)\). 

By Theorem \ref{Frobenius Reciprocity},
\[
\operatorname{Hom}_{P}\!\left(\operatorname{Ind}_{I}^{P} \tau,\; \rho\right)
\cong
\operatorname{Hom}_{I}\!\left(\tau,\; \rho|_I\right).
\]

Then $\tau \hookrightarrow \rho|_I$ corresponds to a non-zero homomorphism
\[
\Phi: \operatorname{Ind}_{I}^{P} \tau \to \rho.
\]


Since \(\rho: P \to \mathrm{GL}(V)\) is irreducible, its only \(P\)-invariant subrepresentations are the zero representation and \(\rho\) itself. 
The \(\operatorname{im}\Phi\) is a \(P\)-invariant subrepresentation of \(\rho\). Hence  \(\operatorname{im}\Phi=V\).


Restrict the induced representation to the normal subgroup \(N\). By Lemma \ref{lem:Mackey's Irreducibility Criterion},
\[
\operatorname{Ind}_{I}^{P} \tau \;\Big|_N \;\cong\; \bigoplus_{gI \,\in\, P/I} (\tau|_N)^g,
\]
where the sum runs over the distinct cosets of \(I\) in \(P\). 

Since \(\tau|_N\) contains \(\mu\), and \(\rho|_N\) is exactly the direct sum of all the distinct conjugates \(\mu^g\)  each appearing with the same multiplicity, the total number of irreducible summands in \(\operatorname{Ind}_{I}^{P} \tau|_N\) is exactly the same as in \(\rho|_N\). 
Thus 
\[
\dim \operatorname{Ind}_{I}^{P} \tau = \dim V.
\]
Then $\Phi$ is an isomorphism.
Therefore,
\[
\operatorname{Ind}_{I}^{P} \tau \cong \rho.
\]

\end{proof}

\begin{lemma}\label{lem:p-group_one-dimensional induced}
Let \(P\) be a finite \(p\)-group. Then every irreducible complex representation of \(P\) is isomorphic to an induced representation
\[
\operatorname{Ind}_{K}^{P} \lambda,
\]
where \(K \le P\) is a subgroup and \(\lambda: K \to \mathbb{C}^{\times}\) is a one-dimensional representation.
\end{lemma}

\begin{proof}
We prove this by induction on the order \(|P|\) of the group \(P\).


Let \((\rho, V)\) be an irreducible complex representation of \(P\). Consider its kernel
\[
K_0 = \ker \rho = \{ g \in P \mid \rho(g) = \mathrm{id}_V \},
\]
which is a normal subgroup of \(P\). The quotient group \(\overline{P} := P/K_0\) is again a finite \(p\)-group. 
Moreover, \[\overline{\rho}:\overline{P} \to\operatorname{Gl}(V),\quad gK_0\mapsto\rho(g)\]is a faithful irreducible representation. 
If we already know the lemma for the smaller group \(\overline{P}\), then there exist a subgroup \(\overline{K} \leq \overline{P}\) and a linear character \(\overline{\lambda} : \overline{K} \to \mathbb{C}^{\times}\) such that
\[
\overline{\rho} \cong \operatorname{Ind}_{\overline{K}}^{\overline{P}} \overline{\lambda}.
\]
Lifting \(\overline{K}\) back to a subgroup \(K \leq P\) and lifting \(\overline{\lambda}\) to a linear character \(\lambda : K \to \mathbb{C}^{\times}\) which is trivial on \(K_0\) gives the desired induced representation for \(\rho\). 

Now suppose \(K_0 = \{1\}\). 


If \(P\) is abelian, then by Lemma \ref{lem:serre_9},  take \(K = P\) and \(\lambda = \rho\). The lemma is proved in this case. 

Suppose that \(P\) is \emph{non-abelian}.



Let \(Z: = Z(P)\) be the center of \(P\), then the quotient group \(P/Z\) is a non-trivial \(p\)-group. Since every non-trivial finite \(p\)-group has a non-trivial center, we can choose an element \(aZ \in Z(P/Z)\) with \(a \notin Z\).
Define $N$ be the abelian group generated by $a$ and elements in $Z$, which is normal in \(P\). 


Actually, if \(N = P\), then \(P/Z\) would be cyclic, which would imply \(P\) is abelian, contradiction. 


Since \(N\) is abelian, by Lemma \ref{lem:serre_9}, all its irreducible complex representations are one-dimensional. Consider the restricted representation \(\rho|_N\). 
Assume that 
\[
\rho|_N \cong \lambda^{\oplus m} \text{  for some positive integer }m,
\]
where \(\lambda : N \to \mathbb{C}^{\times}\) is a one-dimensional representation. Then for every \(n \in N\), \(\rho(n) = \lambda(n) \cdot I_V\).

Since a scalar matrix commutes with every linear operator on \(V\),  \(\rho(n)\) commutes with \(\rho(g)\) for every \(g \in P\). Because \(\rho\) is faithful, \(n\) commutes with every \(g \in P\). Hence \(n \in Z(P)\) for all \(n \in N\). Then \(N \subseteq Z(P)\), contradiction.
Therefore, \(\rho|_N\) must have at least two non-isomorphic irreducible constituents.


Let $(\mu, W )$ be an irreducible subrepresentation of the restricted representation
$(\rho|_N , V )$,and let $I$ be the inertia subgroup of \(\mu\) in \(P\).


By Theorem \ref{Clifford}, there exists an irreducible representation \((\tau,U)\) of \(I\) such that
\[
\rho \cong \operatorname{Ind}_{I}^{P} \tau.
\]
For completeness, Clifford's theorem tells us that \(U|_N\) contains \(W\), and the induced representation from \(I\) to \(P\) recovers the whole of \(V\).

Since \(\rho|_N\) has at least two non-isomorphic constituents, the group \(P\) cannot act transitively on just one isomorphism class. Thus, \(I\) is a proper finite \(p\) -subgroup of \(P\). Apply the induction hypothesis to the irreducible representation \(\tau\) of \(I\). Therefore, there exists a subgroup \(K \leq I\) and a one-dimensional representation \(\lambda : K \to \mathbb{C}^{\times}\) such that
\[
\tau \cong \operatorname{Ind}_{K}^{I} \lambda.
\]
Finally, using the transitivity property of induced representations, we get
\[
\rho \cong \operatorname{Ind}_{I}^{P} \left( \operatorname{Ind}_{K}^{I} \lambda \right)
\cong \operatorname{Ind}_{K}^{P} \lambda.
\]

Thus, every irreducible complex representation of the finite \(p\)-group \(P\) is induced from a one-dimensional representation of some subgroup. 
\end{proof}

\begin{prop}\label{prop:elem_1_dim}
Let \(H = C \times_{grp} P\) be an elementary group, where \(C\) is cyclic of order coprime to \(p\), and \(P\) is a \(p\)-group. Then every irreducible complex representation of \(H\) is isomorphic to
\[
\operatorname{Ind}_{K}^{H} \lambda,
\]
for some subgroup \(K \le H\) and some one-dimensional representation \(\lambda: K \to \mathbb{C}^{\times}\).
\end{prop}

\begin{proof}
Let \((\rho, V)\) be an irreducible complex representation of \(H = C \times_{grp} P\). Since \((|C|, |P|) = 1\), every irreducible representation of a direct product is a tensor product of irreducible representations of the factors. Hence
\[
\rho \cong \alpha \otimes \beta,
\]
where \(\alpha\) is an irreducible representation of \(C\), and \(\beta\) is an irreducible representation of \(P\).

Because \(C\) is cyclic, by Lemma \ref{lem:serre_9}, \(\alpha: C \to \mathbb{C}^{\times}\) is one-dimensional.

By Lemma \ref{lem:p-group_one-dimensional induced}, since \(P\) is a \(p\)-group, there exist a subgroup \(Q \le P\) and a one-dimensional representation \(\lambda_P: Q \to \mathbb{C}^{\times}\) such that
\[
\beta \cong \operatorname{Ind}_{Q}^{P} \lambda_P.
\]

Now set
\[
K := C \times_{grp} Q \le H,
\]
and define a one-dimensional representation \(\lambda: K \to \mathbb{C}^{\times}\) by
\[
\lambda(c, q) := \alpha(c) \cdot \lambda_P(q), \qquad \forall c \in C,\ q \in Q.
\]
Using the standard property of induction for direct products, we have
\[
\operatorname{Ind}_{K}^{H} \lambda
= \operatorname{Ind}_{C \times_{grp} Q}^{C \times_{grp} P} (\alpha \otimes \lambda_P)
= \alpha \otimes \operatorname{Ind}_{Q}^{P} \lambda_P
= \alpha \otimes \beta
\cong \rho.
\]

Therefore \(\rho\) is induced from the one-dimensional representation \(\lambda\) of \(K \le H\), as required.
\end{proof}

\begin{remark}\label{rem:Strong_Brauer}
By Brauer’s Theorem (Theorem \ref{Brauer's Theorem}), every character of irreducible
complex representation of $G$  can be  written as a $\mathbb{Z} $-linear combination of characters induced from elementary subgroups of $G$. By Proposition \ref{prop:elem_1_dim}, every irreducible complex character of an elementary group can be  written as  a character induced from one-dimensional character. Then every character of irreducible
complex representation of $G$ can be  written as a $\mathbb{Z} $-linear combination of characters induced from one-dimensional characters.

Actually, every subgroup of an elementary group is elementary. Then we get the result: every character of irreducible
complex representation of $G$ can be  written as a $\mathbb{Z} $-linear combination of characters induced from one-dimensional characters of elementary subgroups. 
\end{remark}

Now we introduce a special kind of Artin L-function--abelian L-function.

\begin{defn}[Abelian L-Function]
An \textbf{abelian L-function} is an Artin L-function 
\[
L(s, \chi_\rho, L/K)
\]
attached to a representation $(\rho,V)$ of $G$ such that \(\dim V = 1\). 

In this case, the local factor at \(\mathfrak{p}\) simplifies to
\[
L_{\mathfrak{p}}(s, \chi_\rho) = \left(1 - \rho(\mathrm{Frob}_{\mathfrak{p}}) N(\mathfrak{p})^{-s}\right)^{-1}.
\]
\end{defn}
\begin{remark}[Why ``abelian''?]
The name comes from the fact that:

If \(\rho\) is one-dimensional, its image lies in \(\mathbb{C}^{\times}\), which is an abelian group.

Equivalently, such a representation factors through the abelianization \(G^{\mathrm{ab}} = G / [G,G]\).

If \(G\) itself is abelian, then every irreducible representation is one-dimensional, so every Artin L-function for that extension is abelian.

Thus ``abelian L-function'' is a natural shorthand for ``Artin L-function attached to a one-dimensional representation.''
\end{remark}


In the literature, you will often see abelian L-functions referred to as \textbf{Hecke L-functions} (or \textbf{Grossencharakter L-functions}). This is a deep theorem (the \emph{Artin reciprocity law}, part of class field theory), which states that:

\begin{quote}
Every one-dimensional Artin L-function is equal to a Hecke L-function, and conversely.
\end{quote}

For the purposes of this note, we do not need the details of Hecke L-functions. We only need the following analytic fact, which is a classical theorem of Hecke (proved in the 1910s–1920s):


\begin{thm}[Hecke]
Let \(\rho: G \to \mathbb{C}^{\times}\) be a one-dimensional representation. Then its Artin L-function \(L(s, \chi_\rho, L/K)\):
\begin{itemize}
    \item converges absolutely for \(\Re(s) > 1\),
    \item admits a meromorphic continuation to the entire complex plane,
    \item is holomorphic (entire) if \(\rho\) is nontrivial,
    \item has a simple pole at \(s = 1\) if \(\rho\) is trivial.
\end{itemize}
\end{thm}

For a proof, see Hecke's original papers or any standard text on algebraic number theory, such as \cite{neukirch1999algebraic}, \cite{iwaniec2004analytic} and \cite{milne2020cft}.


\begin{thm}[Brauer's Theorem on Artin L-Functions]
Every Artin L-function is meromorphic on the entire complex plane.
\end{thm}

\begin{proof}[Proof Sketch]
By Remark \ref{rem:Strong_Brauer}, \(\chi = \sum_i n_i \operatorname{Ind}_{H_i}^G \psi_i\), where \(n_i \in \mathbb{Z}\) and \(\psi_i\) is a degree-1 character of an elementary subgroup \(H_i\). 

Then
\[
L(s, \chi, L/K) = L(s, \sum_i n_i \operatorname{Ind}_{H_i}^{G} \psi_i, L/K)=\prod_i L(s, \operatorname{Ind}_{H_i}^{G} \psi_i, L/K)^{n_i}= \prod_iL(s, \psi_i, L/L^{H_i}).
\]

Each \(L(s, \psi_i, L/L^{H_i})\) is an abelian L-function, which is known to be holomorphic and has a meromorphic continuation. Hence the product is meromorphic.
\end{proof}

Brauer's theorem only gives meromorphy. The Artin conjecture is a deep open problem in number theory, proven only in special cases (e.g., when \(\rho\) is monomial, or when \(G\) is solvable, by work of Tunnell and Langlands).

\begin{conj}[Artin’s Holomorphy Conjecture]

The \textbf{Artin’s Holomorphy Conjecture } states that for any nontrivial irreducible representation \(\rho\) of \(G\), the associated Artin L-function \(L(s, \chi_\rho, L/K)\) is holomorphic on the entire complex plane (except possibly at \(s = 1\)).
\end{conj}

\subsection{Connection to Modular Forms}

\begin{remark}
For 2-dimensional representations of \(Gal(\overline{\mathbb{Q}}/\mathbb{Q})\), the Artin conjecture is deeply connected to the Langlands program. It is equivalent to the modularity of certain Galois representations, as seen in the proof of Fermat's Last Theorem (Wiles) and the subsequent modularity theorem.
\end{remark}


\nopagebreak
\section*{Conclusion}
As Weyl said, "Symmetry, as wide or as narrow as you may define its meaning, is one idea by which man through the ages has tried to comprehend and create order, beauty, and perfection". 
Representation theory of groups is precisely the language we use to capture, quantify, and employ this symmetry. What this lecture notes lay out is merely the starting point of this long and magnificent road.
